\subsection{Implementation Constraints and Practical Problems}

This section describes the main implementation challenges that arose while implementing the proposed antenna alignment approach and solutions that have been used in terms of software and firmware development. The emphasis is made on the practical constraints observed during simulations and hardware-based implementation.

\subsubsection{RSSI Value Noisiness}

The RSSI values are affected by the phenomena of signal interference, receiver's gain instability, and the presence of thermal noise. Both simulated and real-life scenarios show great variance in repeated RSSI samples taken with the same antenna orientation.

\textbf{Problem:}
\begin{itemize}
    \item Single RSSI sample value led to noisy reward signal
    \item Tiny RSSI value fluctuations led to unwanted actions changes of the RL agent
\end{itemize}

\textbf{Solution applied:}
\begin{itemize}
    \item RSSI values are calculated as an average across sampling period
    \item The direction of RSSI variation (\emph{increase / no change / decrease}) is used to trigger state transition
\end{itemize}

\subsubsection{Errors Related to Quantization and Discretization}

The continuous pan, tilt, and RSSI dimensions were discretized to allow for Q-learning through table lookup.

\textbf{Observed problem:}
\begin{itemize}
\item Very small angle changes could not be captured
\item The optimum angle could potentially fall between two discrete points
\end{itemize}

\textbf{Solution employed:}
\begin{itemize}
\item The discretization interval was chosen carefully based on resolution versus table size considerations
\item The applied controller operated within a realistic search area tested experimentally on physical hardware
\end{itemize}

\subsubsection{Constraints on Memory of Embedded Hardware}

ESP32 has low SRAM and flash memory capacity that makes it difficult to perform dynamic or high dimensional models.

\textbf{Identified problem:}
\begin{itemize}
\item Policies based on neural networks were not suitable for running on ESP32
\item Updating the Q-table in an online fashion would require more complex computation and increase memory requirements
\end{itemize}

\textbf{Proposed solution:}
\begin{itemize}
\item Q-learning was implemented offline on the PC
\item The Q-table from the Q-learning step was only uploaded on ESP32-S3
\item There is no learning done after the deployment phase.
\end{itemize}